% !TeX root = ../papers/nc-p3.tex
\section{二つの相対Picard曲面上での四移動}
\label{sec:finite-residual-picard-surfaces}
\begin{theorem}\label{thm:residual-picard-surface-transport}
定理\ref{thm:finite-quartic-k3-cut}の曲面で $\mathrm{NS}(T)=\mathbb Zh\oplus\mathbb ZF$ とする。
$M_F,M_C$ を二つの鉛筆の、それぞれ次数 $-4,12$ の相対コンパクト化Picard曲面とする。
これらは同型なK3曲面である。共通の曲面 $M$ 上に二つの繊維類 $a,b$ があり
$a^2=b^2=0,\ ab=2$。$a+b$ は豊富かつ大域生成である。
ねじれた普遍層による同値 $U_F:D^b(M,\alpha)\to D^b(T)$ の下で
$\Phi^{-1}[2]$ は $\mathcal O_M(a+b)$ とのテンソル積に対応する。
\end{theorem}
\begin{proof}
$(uh+vF)^2=4u(u+4v)$ なので $T$ に自己交叉 $-2$ の曲線はない。
楕円鉛筆の繊維の各成分は繊維と直交し、Hodge指数定理と随伴により繊維の有理数倍である。
一般繊維が連結で滑らかな完全鉛筆の繊維類は原始的である
（非原始なら原始的なnef等方因子はRiemann--Rochから鉛筆を持ち、倍数系の次元は少なくともその倍数）。
両方の全繊維は従って整である。

ここで用いる相対Picard構成を説明する。整な算術種数一のGorenstein繊維について、次数 $d$ の
階数一無捩れ層は全て安定であり、そのコンパクト化Picardは射影的である
（例えば偏極に関する安定層の射影モジュライのこの繊維部分を取れる）。
基底のエタール開集合上で滑らかな切断 $s$ を取れば、点 $x$ に
$I_{x/\text{繊維}}((d+1)s)$ を対応させる同型で元の曲面に局所同定される。
実際、捩り戻した次数 $-1$ の層は、双対の唯一の切断で繊維の構造層に入り、余核の長さが1である。
この構成は相対双対と切断の基底変換により族で逆となる
（Gorenstein曲線の無捩れ層は反射的であり、次数1の双対の第一コホモロジーは消える）。
よって全空間は滑らかである。座標変更は滑らかな一般繊維の平行移動なので
$R^1\pi_*\mathcal O$ と相対双対線束の基底上の遷移は元の鉛筆と同じである。
いずれも $R^1\pi_*\mathcal O=\mathcal O_{\mathbb P^1}(-2)$、
$\omega_{M/\mathbb P^1}=\pi^*\mathcal O(2)$ となり、全空間はK3である。

普遍層は局所普遍層の基底からの線束のずれをBrauer類に入れて取る。
対応するねじれた導来圏からの変換 $U$ が同値であることは局所でも確認できる。
上の座標では、テンソルとシフトを除き、核は相対対角のイデアルである。
各繊維上で、その核による変換の1移動は
$\mathrm{Cone}(\mathbf R\Gamma(-)\otimes\mathcal O\to\mathrm{Id})$ である。
整な種数一Gorenstein曲線の構造層は完全な1球面対象で、逆は
$\mathrm{Cone}(\mathrm{Id}\to\mathbf R\mathrm{Hom}(-,\mathcal O)^*\otimes\mathcal O)[-1]$ となる。
これは特異繊維でも $H^*(\mathcal O)=k\oplus k[-1]$ と双対化層 $\mathcal O$ のtrace対合により、
二つの評価錐を合成して恒等因子を相殺する球面twistの証明がそのまま適用できる。
相対的にも評価には $\pi^*\mathbf R\pi_*$ を用い、その随伴と相対双対を用いるとこの逆評価が得られる。
平坦な曲線族の完全複体に対する導来基底変換により随伴の単位・余単位は各繊維で同型であり、
導来中山により全空間でも同型である（全空間は正則、相対対角イデアルは両因子に平坦）。
この局所検査は普遍層の線束のずれで変わらない。
また繊維積上の核なので基底の $\mathcal O(1)$ とのテンソル積を保つ。

$V=U_C^{-1}A_0^{-1}U_F$ を取る。定理\ref{thm:octic-fiber-twist-transfer}より
一般点の点層を他方の点層の $-1$ 移動に送る。実際整繊維上で得た階数一無捩れ層はちょうど上の対象である。
この対応は稠密開集合の間の同型である。核の一般平坦性から点層を与える開集合で射となり、
同値ゆえ単射、標数零の同次元曲面なので双有理である。

全点へ延ばすため、両Picard曲面上で、正次元の階数を持たない非零ねじれ層 $Q$ について
$\dim\operatorname{Ext}^1(Q,Q)\ge2$ を示す。$U$ によるMukaiベクトル $x=(r,uh+vF,s)$ の像の記法を
$T$ 側で取ると、$Q$ は一般点層と直交するので $x$ は $(0,F,-4)$ の直交格子に入る。
従って $8u+4r=0$、$x^2=4u(u+4v+s)$ は4の倍数である。
他方 $Q$ の自己Euler標数の負はその一次元サイクルの自乗である。
この値が負なら構成成分の少なくとも一本が負の自乗を持ち、K3の随伴でその曲線は滑らかな有理 $(-2)$ 曲線である。
しかしその曲線上でBrauer類は自明で、その直線束の直像は平方 $-2$ の類を与え、4の倍数性に矛盾する。
ゆえに $x^2\ge0$、Serre双対から所望の不等式が従う。
$C$ 側では $F$ を $C$ に替えて格子が同形であり、直交条件 $8u-12r=0$ は $u=3t,r=2t$ を与え、
再び平方は4の倍数である。この議論は逆Brauer類にも双対化で適用できる。

任意の点の $V$ 像は、対応の稠密開集合のうちその点以外の全点と直交するので、各コホモロジー層は真の閉部分に支えられる。
曲面上の自己Extのスペクトル系列では
$\bigoplus_i\operatorname{Ext}^1(H^i,H^i)$ に入る・出る高次微分はなく、
この次元の和は対象の自己 $\operatorname{Ext}^1$ の次元以下である。
点層でこの上限は2なので、上の不等式からコホモロジー層は一つである。
点の数値類は全点で一定で、一般点と同じ符号付き長さ1の点類だから、それは点層の移動である。
平坦性の開性により移動量は連結な曲面上で一定であり、逆同値にも同じ議論を適用できる。
従って $V$ は曲面の同型のグラフ上のねじれ線束と移動による変換である。
これは各繊維が長さ1の平坦族の台が射のグラフであることからも直接分かる。
従って共役によって相手の鉛筆のテンソル積は引き戻した鉛筆のテンソル積に移る。

二つの繊維類の交叉も符号に依存せず計算できる。$u_0=(0,0,1)$ とすると
$A_0 u_0=(2,-h,1)$。$U_F^{-1}$ は $T$ の一般点層を $M_F$ の滑らかな繊維上の線束の移動へ送る。
これは相対核の随伴を一般点で制限すると通常の楕円曲線上の普遍線束の制限になるためである。
同様に $A_0u_0$ は、今得た同型で $M_C$ の繊維上の線束の移動に対応する。
階数零同士のMukai対合の絶対値が二繊維の交叉数なので $ab=|u_0\cdot A_0u_0|=2$ である。
$a+b$ は基点を持たず正の自乗で、零交叉の曲線があればHodge指数により $(-2)$ 曲線が生じるので豊富である。
最後に命題\ref{prop:quartic-k3-two-pencil-twists}の逆を取り、二つのテンソルに上の同値を施せば結論を得る。
\end{proof}
